\section{Experiments}
\label{sec:experiments}

The experiments ask three questions: how much target-label oracle selection
changes UGDA evaluation, which components explain the source-simulated
development gain, and where the learned covariance correction fails under real
target shifts.
Because the one-time 16-transfer sealed evaluation is not yet complete, all
numbers in this section are explicitly development or source-simulated
results.

\subsection{Experimental setup}

\paragraph{Candidate banks and evaluation scope.}
The four pre-registered development tasks are ACMv9$\rightarrow$Citationv1,
Citationv1$\rightarrow$ACMv9, USA$\rightarrow$BRAZIL, and
BRAZIL$\rightarrow$USA. We report two deliberately separated evaluations.
The equal-protocol real-target comparison uses the pre-registered Gate-1 bank
of 140 frozen candidates per task and is evaluated once by the trusted sealed
process. The source-simulated refinement study uses 675 candidates per task:
five method families (SourceOnly, A2GNN, GRADE, PairAlign, and ADAlign), three
configurations, three seeds, and 15 checkpoints. Its development objective
contains 44 leave-one-source-simulated-shift-out folds and cannot invoke the
sealed evaluator. Candidate trajectories store source-validation outputs,
unlabeled target outputs, embeddings, losses, runtime, and memory metadata.
Frozen implementation values and the source-only sidecar rank gate are listed
in Appendix~\ref{app:implementation}.

\paragraph{Baselines.}
The pre-registered selector suite includes last and best source-validation
selection, target entropy, information maximization, agreement reference,
global committee disagreement, GDE, SND, DEV with source-validation features,
Transfer Score, and ALine-S/ALine-D. All selectors receive the same candidate
bank and public artifacts~\citep{jiang2021disagreement,saito2021tunerightwayunsupervised,
you2019dev,yang2023transferscore,baek2022agreement}. The $0.02278$ reference
reported below is the original frozen SPECTRA selector used to start
controlled source-simulated refinement, not a competing method.
Table~\ref{tab:selector_comparison} supplies the missing equal-protocol
comparison on the four real-target development tasks. The complete
16-transfer sealed comparison remains pending.

\paragraph{Metrics and guardrails.}
For fold $t$, normalized regret is
\begin{equation}
  \nregret_t=
  \frac{R_t(h_{\widehat m})-\min_m R_t(h_m)}
  {\max_m R_t(h_m)-\min_m R_t(h_m)+10^{-12}}.
  \label{eq:nregret}
\end{equation}
Mean normalized regret is primary. We additionally report $\cvar_{20\%}$,
worst-fold regret, median Kendall $\tau$, mean Spearman $\rho$, top-weighted
Kendall, risk-estimation MAE, oracle-F1 gap, top-$5\%$ hit rate, bootstrap
selection stability, localized-gain share, runtime, and protocol audits. A
variant is rejected if it improves the mean by concentrating gains in one fold,
worsens the regret tail, drops Kendall below $0.925$, exceeds 480 seconds, or
accesses target labels.

\subsection{Equal-protocol selector comparison}

\begin{table*}[t]
  \caption{One-time trusted development evaluation on the pre-registered
  Gate-1 real-target bank (four tasks, 140 frozen candidates per task). All
  selectors receive identical public artifacts. This scope is distinct from
  the 675-candidate source-simulated refinement study in
  Table~\ref{tab:core_ablation}. Best values are marked as blue bold
  $\star$; our selector is marked as green bold $\dagger$.}
  \label{tab:selector_comparison}
  \centering
  \small
  \setlength{\tabcolsep}{6pt}
  \resizebox{\textwidth}{!}{%
  \begin{tabular}{lrrrrr}
    \toprule
    Selector & Mean $\nregret$ $\downarrow$ & F1 gap (pt) $\downarrow$ &
    Median $\tau$ $\uparrow$ & Top-$5\%$ hit $\uparrow$ & Micro-F1 (\%) $\uparrow$ \\
    \midrule
    Source validation & 0.2638 & 11.56 & \best{0.761} & 0\% & 53.40 \\
    Last source validation & 0.2388 & 10.93 & 0.683 & 0\% & 54.03 \\
    Entropy & 0.2153 & 8.80 & 0.526 & \best{50\%} & 56.17 \\
    Infomax & 0.1984 & 7.84 & 0.577 & \best{50\%} & 57.12 \\
    DEV & 0.5013 & 27.10 & 0.405 & 0\% & 37.86 \\
    SND & 0.8780 & 47.12 & $-0.582$ & 0\% & 17.84 \\
    Transfer Score & \best{0.1467} & \best{6.33} & 0.356 & \best{50\%} & \best{58.64} \\
    GDE & 0.7539 & 38.73 & 0.001 & 0\% & 26.24 \\
    ALine-S & 0.2638 & 11.56 & \best{0.761} & 0\% & 53.40 \\
    ALine-D & 0.4737 & 26.07 & 0.274 & 0\% & 38.89 \\
    Agreement reference & 0.5365 & 28.81 & $-0.094$ & 0\% & 36.15 \\
    Global disagreement & 0.5414 & 29.14 & $-0.050$ & 0\% & 35.82 \\
    \midrule
    SPECTRA-Static & 0.5414 & 29.14 & $-0.048$ & 0\% & 35.82 \\
    \ours{\method{} (ours)} & \ours{0.2560} & \ours{12.16} &
    \ours{0.695} & \ours{0\%} & \ours{52.81} \\
    \bottomrule
  \end{tabular}%
  }
\end{table*}

Transfer Score is the strongest selector in this real-target development
comparison, attaining mean normalized regret $0.1467$ and selected Micro-F1
$58.64\%$. In contrast, \method obtains $0.2560$ regret and $52.81\%$
Micro-F1, so the present evidence does not support superiority over existing
label-free selectors. The diagnostic is nevertheless informative: \method
reaches median Kendall $\tau=0.695$, substantially above Transfer Score's
$0.356$, but fails to convert this global ranking quality into a competitive
top choice. This is the empirical signature of unreliable correction rather
than an uninformative risk signal: the method orders much of the bank
reasonably, yet its covariance-adjusted top candidate is not conservative
enough under real shift.

\subsection{Source-simulated refinement and core attribution}

Table~\ref{tab:core_ablation} evaluates one controlled change at a time over
the same four development tasks, 44 held-out shift folds, and 675 candidates
per task. The pre-frozen full selector lowers mean normalized regret from the
original selector's $0.02278$ to $0.01415$, a $37.915\%$ relative reduction.
It reaches $\cvar_{20\%}=0.04793$, worst-fold regret $0.08388$, top-weighted
Kendall $0.99462$, and a $95.45\%$ top-$5\%$ hit rate in 327.39 seconds, with
zero target-label accesses and zero protocol violations. This is a
source-simulated refinement result, not a comparison against the real-target
selectors in Table~\ref{tab:selector_comparison}.

\begin{table*}[t]
  \caption{Core source-simulated attribution study (four development tasks,
  44 held-out shifts). ``Cov.'' denotes transported covariance; ``UCB'' is
  bootstrap uncertainty. The frozen full row was fixed before this ablation.
  Lower regret and higher rank/hit metrics are better. Best values are marked
  as blue bold $\star$; the frozen submitted variant is marked as green bold
  $\dagger$; values that are both are orange $\star\dagger$.}
  \label{tab:core_ablation}
  \centering
  \small
  \setlength{\tabcolsep}{7pt}
  \resizebox{\textwidth}{!}{%
  \begin{tabular}{lrrrr}
    \toprule
    Variant & Mean $\nregret$ $\downarrow$ & $\cvar_{20\%}$ $\downarrow$ &
    Median $\tau$ $\uparrow$ & Top-$5\%$ hit $\uparrow$ \\
    \midrule
    Global linear recovery & 0.27767 & 0.62134 & 0.4814 & 0.0\% \\
    Spectral linear recovery & 0.27767 & 0.62134 & 0.4813 & 0.0\% \\
    Spectral + prior + UCB, no cov. & 0.05057 & 0.18877 & 0.8550 & 72.7\% \\
    Global + cov. + UCB (no band conditioning) & 0.01471 & 0.04886 &
    \best{0.9442} & \best{95.5\%} \\
    Spectral + cov. + UCB, no robust refit & \best{0.01397} &
    \best{0.04758} & 0.9398 & \best{95.5\%} \\
    Spectral + cov., no UCB & 0.01452 & 0.04941 & 0.9434 &
    \best{95.5\%} \\
    \ours{\method{} (ours, frozen full)} & \ours{0.01415} &
    \ours{0.04793} & \ours{0.9424} & \oursbest{95.5\%} \\
    \bottomrule
  \end{tabular}%
  }
\end{table*}

The ablation identifies covariance correction, rather than spectral
decomposition by itself, as the main empirical driver. Removing covariance
raises mean regret from $0.01415$ to $0.05057$; the paired mean difference has
a $95\%$ bootstrap interval $[0.01583,0.06026]$. Global and spectral linear
recovery are identical in selection regret, matching the equivalence result in
Appendix~\ref{app:global_equivalence}. Band-conditioned covariance improves the
frozen full row by $3.83\%$ relative to its equal-information global control,
but only four of 44 selections change and the paired interval
$[-0.00073,0.00265]$ crosses zero. The present evidence therefore supports
band-wise heterogeneity, but not a strong claim that spectral conditioning
alone consistently improves top-1 selection.

The solver refinements are secondary. Removing the robust refit slightly
improves mean regret to $0.01397$, while removing UCB changes the frozen mean
from $0.01415$ to $0.01452$; in each comparison at most four folds select a
different candidate. We retain the pre-frozen full row for protocol
consistency, but treat robust fitting and bootstrap uncertainty as
implementation details rather than contributions.

\subsection{Do spectral bands expose different error regimes?}

The source-only preflight evaluates 176 leave-one-simulated-shift-out
comparisons over the planned 16-task calibration bank, with 675 candidates per
task. Table~\ref{tab:covariance_band} rejects a simplistic explanation that
every graph frequency decorrelates model errors. The low-frequency band has
lower absolute error correlation than the global statistic in $98.9\%$ of
task-shift comparisons, but the middle and high bands are more correlated on
average. Conversely, the middle band has the lowest covariance-transport NMAE.
The supported mechanism is therefore \emph{band-wise heterogeneity}: different
frequencies expose different correlation and transportability regimes that a
band-conditioned recovery rule can treat differently.

\begin{table}[t]
  \caption{Source-simulated covariance diagnostics. ``Below global'' is the
  fraction of task-shift comparisons whose band correlation is below the
  corresponding global correlation.}
  \label{tab:covariance_band}
  \centering
  \small
  \begin{tabular}{lrrrr}
    \toprule
    Statistic & Global & Low & Mid & High \\
    \midrule
    Mean $|$error correlation$|$ & 0.487 & \best{0.381} & 0.573 & 0.617 \\
    Transport NMAE & 0.092 & 0.227 & \best{0.069} & 0.186 \\
    Below global & -- & \best{98.9\%} & 0.0\% & 5.1\% \\
    \bottomrule
  \end{tabular}
\end{table}

\subsection{Failure analysis and scope}

The dominant unresolved issue is cross-family covariance transport. A
low-rank risk-correction variant (SARC) performs poorly under shift-family
holdout, with mean regret $0.06977$; moreover, descriptor distance has
Spearman correlation $-0.04954$ with risk-correction distance. Thus the
current unlabeled descriptor geometry does not reliably identify which source
simulation has a transferable error-covariance correction. Full covariance
interactions remain selection-relevant, but predicting them for an unseen
shift family is not solved by the current model.

The source-simulated results support covariance-aware risk recovery and expose
band-wise covariance heterogeneity, whereas the real-target comparison shows
that the operational selector is not yet competitive with Transfer Score.
Together, these results indicate that the hard problem is not recovering risk
from disagreement once the correction is known, but deciding when a
source-simulated covariance correction should be trusted on an unseen target.
The evidence therefore does not establish full UGDA state of the art, broad
method-ranking reversals, or generalization to all 16 ADAlign transfers. Those
claims require a frozen one-time target evaluation across Citation, Airport,
Blog, and Twitch families, paired task-level uncertainty, equal-information
comparison to global calibration, and explicit reliability tests such as
leave-one-shift-family-out covariance correction. We preserve this boundary
throughout the paper.
